% SPDX-License-Identifier: CC-BY-SA-4.0
% Copyright (c) 2025 Luis Alonso
%
% This file is part of luis_alonso/Notas-Japones.
%
% Licensed under Creative Commons Attribution-ShareAlike 4.0 International (CC BY-SA 4.0).
% You may share and adapt this material for any purpose,
% as long as you provide attribution and distribute adaptations under the same license.
%
% License: https://creativecommons.org/licenses/by-sa/4.0/
% Source:  https://git.lalonso.com/luis_alonso/Notas-Japones
%
% Note: Third-party content (if any) is excluded from this license and is marked where used.

\documentclass[a4paper,lualatex,ja=standard,base=10pt]{bxjsarticle}

\title{\ruby{漢|字}{かん|じ} LV4}
\author{アロンゾ　ルイス}

\setmainjfont{Noto Serif CJK JP}
\usepackage{luatexja-fontspec}
\usepackage{kanjicard} % Custom Kanji card
\usepackage{fancyhdr}

% --- Page style ---
\pagestyle{fancy}
\fancyhf{} % clear header/footer

% Footer: name + copyright on left, page number on right
\fancyfoot[L]{Luis Alonso\ \copyright\ \the\year}
\fancyfoot[R]{\thepage}

% Optional: remove header rule and keep a footer rule
\renewcommand{\headrulewidth}{0pt}
\renewcommand{\footrulewidth}{0.4pt}

% Increase space reserved for header area
\setlength{\headheight}{0pt}
\setlength{\headsep}{6pt} % distance from header to text

\fancypagestyle{plain}{%
  \fancyhf{}
  \fancyfoot[L]{Luis Alonso\ \copyright\ \the\year}
  \fancyfoot[R]{\thepage}
  \renewcommand{\headrulewidth}{0pt}
  \renewcommand{\footrulewidth}{0.4pt}
}

\begin{document}
\setcounter{page}{1} % start page number here

\raggedbottom
\maketitle{}
\thispagestyle{fancy}

\begin{kanjicard}{私}{わたし}{シ}
  \KExample{私}{わたし}{Yo}
  \KExample{私立大学}{し・りつ・だい・がく}{Universidad privada}
\end{kanjicard}

\begin{kanjicard}{父}{ちち}{フ/ブ/ホ}
  \KExample{父}{ちち}{Padre (mio)}
  \KExample{祖父}{そ・ふ}{Abuelo (mio)}
  \KExample{お父さん}{お・とう・さん}{Padre de alguien mas}
\end{kanjicard}

\begin{kanjicard}{母}{はは}{ボ/ボウ/モ}
  \KExample{母}{はは}{Madre (mia)}
  \KExample{祖母}{そ・ぼ}{Abuela (mia)}
  \KExample{お母さん}{お・か・さん}{Madre (de alguien mas)}
\end{kanjicard}

\begin{kanjicard}{子}{こ}{シ/ス}
  \KExample{子ども}{こ・ども}{Niño/a}
  \KExample{息子}{むす・こ}{Hijo}
\end{kanjicard}

\begin{kanjicard}{国}{くに}{コク}
  \KExample{国}{くに}{Pais}
  \KExample{外国}{がいこく}{Pais extranjero}
  \KExample{中国}{ちゅうごく}{China}
\end{kanjicard}

\begin{kanjicard}{男}{おとこ}{ダン}
  \KExample{男}{おとこ}{Masculino}
  \KExample{男の子}{おとこ・の・こ}{Niño}
  \KExample{男の人}{おとこ・の・ひと}{Hombre}
\end{kanjicard}

\begin{kanjicard}{女}{おんな}{ショ}
  \KExample{女}{おんな}{Mujer/Femenino}
  \KExample{女の子}{おんな・の・こ}{Niña}
  \KExample{女の人}{おんな・の・ひと}{Mujer}
\end{kanjicard}

\begin{kanjicard}{人}{ひと}{ジン/ニン}
  \KExample{人}{ひと}{Persona}
  \KExample{メキシコ人}{メキシコ・じん}{Nacionalidad mexicana}
  \KExample{大人}{おとな}{Adulto}
  \KExample{友人}{ゆう・じん}{Amigo}
\end{kanjicard}

\begin{kanjicard}{外}{そと}{ガイ}
  \KExample{海外}{かいがい}{Extranjero}
  \KExample{外}{そと}{Exterior/Fuera}
  \KExample{外国人}{がいこくじん}{Persona extranjera}
\end{kanjicard}

\begin{kanjicard}{語}{かたります}{ゴ}
  \KExample{言語}{げんご}{Idioma}
  \KExample{外国語}{がいこくご}{Idioma extranjero}
  \KExample{日本語}{にほんご}{Idioma Japonés}
\end{kanjicard}

\begin{kanjicard}{本}{}{ほん}
  \KExample{本}{ほん}{Libro}
  \KExample{日本}{にほん}{Japón}
  \KExample{本屋}{ほんや}{Librería}
\end{kanjicard}

\begin{kanjicard}{英}{}{エイ}
  \KExample{英語}{えいご}{Idioma Inglés}
  \KExample{英国}{えいこく}{Gran Bretaña}
\end{kanjicard}

\begin{kanjicard}{中}{なか}{ジュウ/チュウ}
  \KExample{中国語}{ちゅう・ごく・ご}{Idioma Chino}
  \KExample{中}{なか}{Dentro}
  \KExample{勉強中}{べん・きょう・ちゅう}{Estudiando actualmente}
\end{kanjicard}

\begin{kanjicard}{人}{ひと}{ジン/ニン}
  \KExample{人}{ひと}{Persona}
  \KExample{中国人}{ちゅうごくじん}{Nacionalidad china}
\end{kanjicard}

\begin{kanjicard}{好}{好き}{コウ}
  \KExample{好き}{すき}{Gustar}
  \KExample{友好的}{ゆうこうてき}{Amigable}
\end{kanjicard}

\begin{kanjicard}{読}{よ・みます}{ドク/トク}
  \KExample{読書}{どく・しょ}{Lectura}
  \KExample{読みます}{よ・みます}{Leer}
  \KExample{読解}{どっ・かい}{Comprension lectora}
\end{kanjicard}

\begin{kanjicard}{書}{か・きます}{ジョ}
  \KExample{読書}{どく・しょ}{Lectura}
  \KExample{書きます}{か・きます}{Escribir}
  \KExample{書道}{しょ・どう}{Caligrafía}
\end{kanjicard}

\begin{kanjicard}{何}{なに/なん}{}
  \KExample{何}{なに}{Qué?}
  \KExample{何人}{なん・にん}{Cuantas personas?}
\end{kanjicard}

\begin{kanjicard}{春}{はる}{シュン}
  \KExample{春}{はる}{Primavera}
  \KExample{青春}{せい・しゅん}{Juventud}
  \KExample{春休み}{はる・やす・み}{Vacaciones de primavera}
\end{kanjicard}

\begin{kanjicard}{夏}{なつ}{}
  \KExample{夏}{なつ}{Verano}
  \KExample{夏休み}{なつ・やす・み}{Vacaciones de verano}
\end{kanjicard}

\begin{kanjicard}{秋}{あき}{シュウ}
  \KExample{秋}{あき}{Otoño}
  \KExample{秋分}{しゅう・ぶん}{Equinoccio de otoño}
  \KExample{秋祭り}{あき・まつ・り}{Festival de otoño}
\end{kanjicard}

\begin{kanjicard}{冬}{ふゆ}{トウ}
  \KExample{冬}{ふゆ}{Invierno}
  \KExample{冬休み}{ふゆ・やす・み}{Vacaciones de invierno}
\end{kanjicard}

\begin{kanjicard}{今}{いま}{コン}
  \KExample{今}{いま}{Ahora}
  \KExample{今春}{こん・しゅん}{Esta primavera}
  \KExample{今秋}{こん・しゅう}{Este otoño}
  \KExample{今冬}{こん・とう}{Este invierno}
  \KExample{今夏}{こん・か}{Este verano}
  \KExample{今日}{きょう}{Hoy}
\end{kanjicard}

\begin{kanjicard}{花}{はな}{カ/ケ}
  \KExample{花}{はな}{Flor}
  \KExample{花火}{はな・び}{Fuegos artificiales}
  \KExample{花見}{はな・み}{Observacion de flores}
  \KExample{花瓶}{か・びん}{Florero}
\end{kanjicard}

\begin{kanjicard}{海}{うみ}{カイ}
  \KExample{海}{うみ}{Mar}
  \KExample{海外}{かい・がい}{Extranjero}
\end{kanjicard}

\begin{kanjicard}{山}{やま}{サン}
  \KExample{山}{やま}{Montaña}
  \KExample{山登り}{やま・のぼ・り}{Montañismo}
  \KExample{富士山}{ふ・じ・さん}{Monte Fuji}
\end{kanjicard}

\begin{kanjicard}{川}{かわ}{}
  \KExample{川}{かわ}{Rio}
  \KExample{ナイル川}{ナイル・がわ}{Río Nilo}
\end{kanjicard}

\begin{kanjicard}{日}{ひ/か}{ジッ/ニチ}
  \KExample{一日中}{いち・にち・じゅう}{Todo el día}
  \KExample{明日}{あした}{Mañana}
\end{kanjicard}

\begin{kanjicard}{天}{}{テン}
  \KExample{天気予報}{てん・き・よ・ほう}{Pronostico del clima}
  \KExample{天使}{てん・し}{Ángel}
\end{kanjicard}

\begin{kanjicard}{気}{いき}{キ}
  \KExample{天気}{てん・き}{Clima}
  \KExample{人気}{にん・き}{Popularidad}
  \KExample{気持ち}{き・も・ち}{Sentimiento}
  \KExample{病気}{びょう・き}{Enfermedad}
  \KExample{電気}{でん・き}{Electricidad}
\end{kanjicard}

\begin{kanjicard}{晴}{は・れ}{セイ}
  \KExample{晴れ}{は・れ}{Despejado (clima)}
  \KExample{素晴らしい}{す・ば・らしい}{Esplendido, maravilloso}
\end{kanjicard}

\begin{kanjicard}{雨}{あめ}{ウ}
  \KExample{雨}{あめ}{Lluvia}
  \KExample{梅雨}{つ・ゆ}{Temporada de lluvias}
  \KExample{大雨}{おお・あめ}{Chubasco, lluvia pesada}
\end{kanjicard}

\begin{kanjicard}{雪}{ゆき}{セツ}
  \KExample{雪}{ゆき}{Nieve}
  \KExample{白雪}{しら・ゆき}{Nieve blanca}
  \KExample{雪白}{せっ・ぱく}{Puro, inmaculado}
\end{kanjicard}

\begin{kanjicard}{雲}{くも}{ウン}
  \KExample{雲}{くも}{Nube}
  \KExample{雨雲}{あま・ぐも}{Nubes de lluvia}
\end{kanjicard}

\begin{kanjicard}{風}{かぜ}{ふ/フウ}
  \KExample{風}{かぜ}{Viento}
  \KExample{台風}{たい・ふう}{Huracán, tifón}
\end{kanjicard}

\begin{kanjicard}{空}{そら}{クウ}
  \KExample{空}{そら}{Cielo}
  \KExample{空気}{くう・き}{Aire, atmosfera}
  \KExample{空港}{くう・こう}{Aeropuerto}
\end{kanjicard}

\begin{kanjicard}{町}{まち}{チョウ/テイ}
  \KExample{町}{まち}{Ciudad}
  \KExample{町家}{まち・や}{Casa tradicional}
\end{kanjicard}

\begin{kanjicard}{店}{みせ}{テン}
  \KExample{店}{みせ}{Tienda}
  \KExample{喫茶店}{きっ・さ・てん}{Cafetería}
  \KExample{店員}{てん・いん}{Empleado de tienda}
\end{kanjicard}

\begin{kanjicard}{多}{おお・い}{タ}
  \KExample{多い}{おお・い}{Muchos}
  \KExample{多分}{た・ぶん}{Tal vez}
\end{kanjicard}

\begin{kanjicard}{少}{すく・ない}{ショウ}
  \KExample{少ない}{すく・ない}{Poco, poca cantidad}
  \KExample{少し}{すこ・し}{Poco, insuficiente}
\end{kanjicard}

\begin{kanjicard}{高}{たか・い}{コウ}
  \KExample{高い}{たか・い}{Alto, caro}
  \KExample{高校}{こう・こう}{Preparatoria}
  \KExample{最高}{さい・こう}{Lo mejor, mas alto}
\end{kanjicard}

\begin{kanjicard}{安}{やす・い}{アン}
  \KExample{安い}{やす・い}{Barato}
  \KExample{安心}{あん・しん}{Tranquilidad de espíritu}
  \KExample{安全}{あん・ぜん}{Seguridad}
\end{kanjicard}

\begin{kanjicard}{広}{ひろ・い}{コウ}
  \KExample{広い}{ひろ・い}{Amplio}
  \KExample{広場}{ひろ・ば}{Plaza}
  \KExample{広告}{こう・こく}{Anuncio publicitario}
\end{kanjicard}

\begin{kanjicard}{道}{みち}{ドウ/トウ}
  \KExample{道}{みち}{Camino}
  \KExample{書道}{しょ・どう}{Caligrafía}
  \KExample{近道}{ちか・みち}{Atajo}
\end{kanjicard}

\begin{kanjicard}{通}{とお・り}{ツウ}
  \KExample{通り}{とお・り}{Calle}
  \KExample{普通}{ふ・つう}{Ordinario}
\end{kanjicard}

\begin{kanjicard}{右}{みぎ}{ユウ}
  \KExample{右}{みぎ}{Derecha}
  \KExample{右手}{みぎ・て}{Diestro}
  \KExample{右側}{みぎ・がわ}{Lado derecho}
\end{kanjicard}

\begin{kanjicard}{左}{ひだり}{サ}
  \KExample{左}{ひだり}{Izquierda}
  \KExample{左手}{ひだり・て}{Zurdo}
  \KExample{左側}{ひだり・がわ}{Lado izquierdo}
\end{kanjicard}

\begin{kanjicard}{一}{ひと・つ}{イチ}
  \KExample{一つ}{ひと・つ}{Una cosa}
  \KExample{一人}{ひとり}{Una persona}
\end{kanjicard}

\begin{kanjicard}{二}{ふた・つ}{ニ}
  \KExample{二つ}{ふた・つ}{Dos cosas}
  \KExample{二人}{ふたり}{Dos personas}
  \KExample{二度と}{に・ど・と}{Nunca más}
\end{kanjicard}

\begin{kanjicard}{赤}{あか}{セキ}
  \KExample{赤い}{あか・い}{Rojo}
  \KExample{赤ちゃん}{あか・ちゃん}{Bebé}
  \KExample{赤点}{あか・てん}{Marca reprobatoria}
\end{kanjicard}

\begin{kanjicard}{青}{あお}{セイ/ショウ}
  \KExample{青い}{あお・い}{Azul}
  \KExample{青空}{あお・ぞら}{Cielo azul}
  \KExample{青春}{せい・しゅん}{Juventud, adolescencia}
\end{kanjicard}

\begin{kanjicard}{黒}{くろ}{クロ}
  \KExample{黒い}{くろ・い}{Negro}
  \KExample{黒板}{こく・ばん}{Pizarrón (gis)}
  \KExample{黒雲}{くろ・くも}{Nubes negras}
\end{kanjicard}

\begin{kanjicard}{白}{しろ}{ハク}
  \KExample{白い}{しろ・い}{Blanco}
  \KExample{白板}{はく・ばん}{Pizarrón (pintarron)}
  \KExample{告白}{こく・はく}{Confesión (amor)}
\end{kanjicard}

\end{document}

% 沢山 es un adverbio va antes de lo que describes 多い es un adjetivo
% normalmente va despues
% あたり en esta zona (en esta zona hay ...)
